\section{Model Evaluation and Extensions}

\subsection{Strengths}

\begin{enumerate}
    \item \textbf{Hierarchical forecasting preserves heterogeneity and consistency:}
    Problem 1 recovers 18 cross-classified demands from three forecast levels. Validation and test WAPE are 0.7241 and 0.6519, and all 84 rolling windows outperform the 24-hour same-hour baseline.
	
    \item \textbf{Multiobjective scheduling exposes practical trade-offs:}
    Problem 2 uses fixed normalization, and full-horizon recomputation confirms that all 50000 jobs satisfy hard constraints. Cost falls by 507.34 units of CNY $10^4$ and renewable utilization rises by 0.223 percentage points, with the increased emissions and latency also reported.
	
    \item \textbf{The storage model retains the specified load convention:}
    Problem 3 fixes the specified loads and optimizes only energy decisions. Relative to no storage, cost falls by 3936.19 units of CNY $10^4$, with lower emissions, peaks, and fluctuations. SOC and terminal constraints pass verification.
	
    \item \textbf{Joint scheduling captures computing--storage--electricity interactions:}
    Problem 4 recomputes sequential and joint job schedules with consistent energy accounting. Joint jobs lower cost by 65.67 units of CNY $10^4$, mean latency by 2.9851 ms, and renewable nonutilization by 0.0258 percentage points, while raising emissions and peaks by 41.08 tCO$_2$ and 30.61 MW. This fully reports the energy-side cost of service-quality improvements. Both schedules pass all 67 independent checks.
\end{enumerate}


\subsection{Limitations}

\begin{enumerate}
    \item \textbf{Forecasts depend on stable recent demand structure:}
    Recovery of the 18 demands relies on historical structure, and its accuracy change relative to direct forecasting is not statistically significant ($p=0.2163$). Abrupt business changes require dynamic windows or time-varying weights.
	
    \item \textbf{The Problem 2 rolling heuristic does not guarantee global optimality:}
    The maximum normalized-objective deviation from exact models in difficult windows is 293.40\%. Resource-constrained neighborhoods could be expanded and critical windows given longer solve times.
	
    \item \textbf{Extreme scenarios and uncertainty receive limited treatment:}
    Problems 3 and 4 use deterministic prices, carbon intensity, and renewable sequences, excluding forecasting errors and storage degradation. Progressively tightened constraints and robust or stochastic optimization could be introduced.
\end{enumerate}


\subsection{Extensions}

The framework can support multiregional data centers, computing networks, and campus source--load--storage scheduling, but must be recalibrated to actual job windows, networks, capacities, PUE, prices, carbon intensity, renewables, storage, and SLAs. Introducing preemption, interregional electricity exchange, or forecasting errors requires additional state and uncertainty constraints and renewed validation. The transferable element is the modeling framework, not the parameters or numerical results in this paper.
